\section*{Capítulo \ref{ch:b1:induction} --- Indução eletromagnética e aplicações}

\begin{solution}{exo:b1:induction:1}
$\Phi = NBS\cos30^\circ = 100 \times 0.2 \times 10^{-3} \times 0.866 = \qty{1.7e-2}{Wb}$;
$\langle e\rangle = \Delta\Phi/\Delta t = \qty{1.7}{V}$ em \qty{10}{ms}, e \qty{17}{V} em
\qty{1}{ms}. A corrente circula de modo a recriar o fluxo que desaparece: o seu
campo tem o sentido do antigo $\vect B$.
\end{solution}

\begin{solution}{exo:b1:induction:2}
$e = B\ell v = 0.5 \times 0.2 \times 2 = \qty{0.20}{V}$; $i = \qty{2.0}{A}$; $F = B\ell i =
\qty{0.20}{N}$ freando, de modo que quem empurra aplica \qty{0.20}{N}; $Fv = \qty{0.40}{W}
= Ri^2$.
\end{solution}

\begin{solution}{exo:b1:induction:3}
$L = \mu_0n^2S\ell = 4\pi \times 10^{-7} \times 10^6 \times 4 \times 10^{-4} \times 0.2 =
\qty{0.10}{mH}$; $W = \tfrac12LI^2 = \qty{1.3}{mJ}$; $\tau = L/R = \qty{0.1}{ms}$.
$B = \mu_0nI = \qty{6.3}{mT}$, $B^2/2\mu_0 = \qty{15.7}{J/m^3}$ sobre
$S\ell = \qty{8e-5}{m^3}$: \qty{1.3}{mJ}.
\end{solution}

\begin{solution}{exo:b1:induction:4}
$N_2/N_1 = 12/230 = 0.052$; $N_1 = 50/0.052 = 960$ espiras; $i_1 = 2 \times 0.052 =
\qty{0.10}{A}$. Uma corrente contínua faz um fluxo constante: sem $\dd\varphi/\dd t$, sem
\omterm{def:b1:dc-circuits:voltage}{tensão} no secundário --- e o primário, um curto, queimaria.
\end{solution}

\begin{solution}{exo:b1:induction:5}
(a) O fluxo (no sentido da aproximação) cresce: a corrente cria um campo
que se opõe a ele --- o seu próprio polo norte fica de frente para o norte do ímã --- e o
anel é empurrado para longe. (b) Uma corrente cujo campo se opõe ao crescimento; o
anel, um dipolo antialinhado com $\vect B$, é empurrado para o campo mais
fraco. (c) Ao cair, o fluxo cresce: mesma corrente que em (a), e o anel é
freado. (d) Cada trecho de tubo à frente do ímã vê o fluxo crescer, e
atrás dele, diminuir; as duas correntes induzidas freiam o ímã; na
velocidade em que a \omterm{thm:b1:magnetostatics:laplace}{força de Laplace} equilibra o peso, o ímã cai
uniformemente e a sua \omterm{def:b1:work-and-energy:conservative}{energia potencial} perdida vira calor Joule no tubo.
\end{solution}

\begin{solution}{exo:b1:induction:6}
$E_{\max} = NBS\omega = 200 \times 2 \times 10^{-3} \times 0.1 \times 314 = \qty{12.6}{V}$, eficaz
\qty{8.9}{V}; $I_{\max} = \qty{1.26}{A}$, $\Gamma_{\max} = NBSI_{\max} = \qty{0.050}{N.m}$
(quando $\sin\omega t = 1$); $\langle P\rangle = E_{\max}^2/2R = \qty{7.9}{W}$ $=
\langle\Gamma\omega\rangle = \tfrac12 \times 0.05 \times 314$; quando $e = 0$ a corrente é
nula e o \omterm{def:b1:angular-momentum:moment}{torque} se anula --- o \omterm{def:b1:angular-momentum:moment}{torque} pulsa a \qty{100}{Hz}.
\end{solution}

\begin{solution}{exo:b1:induction:7}
$e = Bav$, $i = Bav/R$, $F = B^2a^2v/R = 10\,v$ (N). $m\,\dd v/\dd t = -B^2a^2v/R$:
$v = v_0\eu^{-t/\tau}$, $\tau = mR/B^2a^2 = \qty{0.10}{s}$; distância $v_0\tau =
\qty{10}{cm}$ (se a região de campo for longa o bastante). A \omterm{thm:b1:work-and-energy:ke}{energia cinética},
\qty{0.5}{J}, é dissipada como $Ri^2$ na espira.
\end{solution}

\begin{solution}{exo:b1:induction:8}
Fluxo do campo do solenoide $\mu_0n_1i_1$ através das $N_2$ espiras: $M =
\mu_0n_1N_2S_2 = 4\pi \times 10^{-7} \times 2000 \times 50 \times 2 \times 10^{-4} = \qty{25}{\micro H}$.
Rampa: $e = M\,\dd i_1/\dd t = \qty{2.5}{mV}$; a \qty{10}{kHz} e \qty{1}{A}: $e_{\max} =
M\omega I = 25 \times 10^{-6} \times 6.3 \times 10^4 = \qty{1.6}{V}$. $M$ é simétrica:
o mesmo número dá o fluxo que a bobina pequena envia através do
solenoide --- muito mais difícil de calcular diretamente.
\end{solution}

\begin{solution}{exo:b1:induction:9}
$W = B^2V/2\mu_0 = 2.25/(2.51 \times 10^{-6}) = \qty{0.90}{MJ}$; $L = 2W/I^2 =
\qty{7.2}{H}$; $h = W/mg = \qty{92}{m}$; hélio: $9 \times 10^5/2600 = \qty{350}{L}$.
Campo da Terra: $(5 \times 10^{-5})^2/2\mu_0 = \qty{1e-3}{J/m^3}$, \qty{50}{mJ} na
sala.
\end{solution}

\begin{solution}{exo:b1:induction:10}
(a) $E - B\ell v = Ri$ (a força contraeletromotriz se opõe à fonte), $m\,\dd v/\dd t =
B\ell i$: $m\,\dd v/\dd t = B\ell(E - B\ell v)/R$. (b) $v_\infty = E/B\ell$, $\tau =
mR/B^2\ell^2$. (c) Multiplique a equação elétrica por $i$: $Ei = Ri^2 +
B\ell vi$, e $B\ell vi = Fv$ é a potência mecânica; \omterm{def:b1:heat-engines:efficiency}{rendimento} $Fv/Ei =
B\ell v/E = v/v_\infty$ --- $100\%$ só a vazio, quando nada é
entregue. (d) $v_\infty = 2/(0.5 \times 0.2) = \qty{20}{m/s}$, $\tau = 0.05 \times
0.1/0.01 = \qty{0.5}{s}$. (e) $F = B\ell I = 1 \times 0.1 \times 10^6 = \qty{1e5}{N}$
(tomando $\ell = \qty{10}{cm}$), $a = \qty{1e7}{m/s^2}$, $v = \sqrt{2a \times 5} =
\qty{1e4}{m/s}$ --- dez quilômetros por segundo, a promessa e o
problema dos canhões eletromagnéticos (os trilhos se erodem).
\end{solution}

\begin{solution}{exo:b1:induction:11}
(a) No raio $r$ o metal se move a $v = \omega r$; $\vect v\wedge\vect B$ é
radial, de módulo $\omega rB$: $e = \int_0^R\omega Br\,\dd r = \tfrac12B\omega R^2$.
(b) $\omega = 314$: $e = 0.5 \times 1 \times 314 \times 0.01 = \qty{1.6}{V}$. (c) Uma
só ``espira'', logo alguns volts no máximo; em \qty{1}{m\ohm}: \qty{1.6}{kA}, \omterm{def:b1:angular-momentum:moment}{torque}
$= P/\omega = ei/\omega = 2500/314 = \qty{8}{N.m}$. (d) O circuito escova--raio--
borda--escova varre a área $\tfrac12R^2\omega\,\dd t$ em cada $\dd t$: $\dd\Phi/\dd t =
\tfrac12BR^2\omega$ --- é o fluxo através do circuito móvel e deformável.
\end{solution}

\begin{solution}{exo:b1:induction:12}
(a) $e_2 = -M\,\dd i_1/\dd t$, de pico $M\omega I_1 = 10^{-6} \times 1.57 \times 10^5 \times 30 =
\qty{4.7}{V}$. (b) $I_2 = 4.7/0.01 = \qty{470}{A}$ de pico; $\langle P\rangle = R_pI_2^2/2
= \qty{1.1}{kW}$. (c) A fem é $\propto\omega$: uma frequência alta dá uma fem
útil com um $M$ modesto; o vidro é isolante, e nada circula nele
--- o calor nasce dentro do metal. (d) A corrente da panela induz na
bobina $M\omega I_2 = 10^{-6} \times 1.57 \times 10^5 \times 470 = \qty{74}{V}$ de pico,
em fase com $i_1$ (a corrente da panela não se atrasa em relação a $e_2$, sendo
resistiva, e $e_2$ se atrasa de \qty{90}{\degree} em relação a $i_1$\ldots\ de modo que a
fem refletida vale $-M\,\dd i_2/\dd t$, em oposição de fase com a \omterm{def:b1:dc-circuits:voltage}{tensão}
\emph{motriz} de $i_1$): o gerador fornece $\tfrac12 \times 74 \times 30 =
\qty{1.1}{kW}$ a mais --- exatamente o que a panela dissipa.
\end{solution}

\begin{solution}{pb:b1:induction:1}
\textbf{1.} $\Phi = NBS\cos\omega t$, $e = NBS\omega\sin\omega t$; $NBS = 300 \times
0.3 \times 3 \times 10^{-4} = \qty{0.027}{Wb}$; $\omega = v/r = 5/0.01 = \qty{500}{rad/s}$:
\omterm{def:b1:wave-propagation:signal}{amplitude} \qty{13.5}{V}, frequência $\omega/2\pi = \qty{80}{Hz}$.

\textbf{2.} $I_{\text{rms}} = 13.5/\sqrt2/16 = \qty{0.60}{A}$; lâmpada $RI^2 = \qty{4.3}{W}$
sob \qty{7.2}{V}: uma lâmpada de \qty{6}{V} e \qty{3}{W} fica sobrecarregada e não
dura.

\textbf{3.} Mecânica $=$ elétrica: $(R + r_c)I_{\text{rms}}^2 = 16 \times 0.36 =
\qty{5.7}{W}$; $F = P/v = 5.7/5 = \qty{1.1}{N}$ --- metade do atrito de rolamento;
o ciclista sente.

\textbf{4.} O binário de Laplace da corrente induzida se opõe à rotação
(Lenz); a energia da lâmpada vem das pernas do ciclista, \qty{5.7}{W} dela.

\textbf{5.} $\underline Z = R + r_c + jL\omega$; $I = NBS\omega/\sqrt{(R + r_c)^2 + L^2\omega^2}$;
para $L\omega \gg R + r_c$, $I \to NBS/L = 0.027/0.02 = \qty{1.35}{A}$. $70\%$: $L\omega
= 0.7(R + r_c)/\sqrt{1 - 0.49} = 15.7$, $\omega = \qty{784}{rad/s}$, $v = \qty{7.8}{m/s}
= \qty{28}{km/h}$.

\textbf{6.} Sem $L$: $e_{\max} = 3 \times 13.5 = \qty{40.5}{V}$, $I_{\text{rms}} =
\qty{1.8}{A}$, \qty{39}{W} na lâmpada --- ela queima. Com $L$: $I = 40.5/\sqrt{256 +
900} = \qty{1.2}{A}$ de pico, \qty{8.5}{W}: quente, mas dez vezes menos.

\textbf{7.} \qty{9}{km/h}: $e_{\max} = 6.75$, $|Z| = \sqrt{256 + 25} = 16.8$, $I = 0.40$,
$P = \tfrac12 \times 0.16 \times 12 = \qty{1.0}{W}$: fraca. \qty{36}{km/h}: $e_{\max} = 27$,
$|Z| = 25.6$, $I = 1.05$, $P = \qty{6.6}{W}$. Quatro vezes a velocidade, $2.6$ vezes
a corrente: a indutância regula.

\textbf{8.} Só o \omterm{def:b1:non-inertial-frames:frames}{movimento relativo} importa: o fluxo através da bobina fixa
varia do mesmo modo, mesma fem. Não são precisos contatos deslizantes
(escovas) para tirar a corrente de uma bobina que gira.

\textbf{9.} Sem movimento, sem variação de fluxo, sem luz --- daí o \omterm{def:b1:potential-capacitors:capacitance}{capacitor} ou
a pequena bateria (``luz de parada'') carregada durante o pedal.

\textbf{10.} Cada elemento de fio é perpendicular ao $\vect B$ radial:
$F = B\ell i$, axial. Movendo-se a $v$, cada elemento vê $\vect v\wedge\vect B$
ao longo do fio: $e = -B\ell v$, que se opõe à corrente que produz o
movimento (\omterm{thm:b1:dc-circuits:kirchhoff}{convenção receptor}: uma \omterm{def:b1:dc-circuits:voltage}{tensão} $B\ell v$).

\textbf{11.} $u = Ri + L\,\dd i/\dd t + B\ell v$; $m\,\dd v/\dd t = -kx - hv + B\ell i$.

\textbf{12.} $\underline i = (\underline u - B\ell\underline v)/(R + jL\omega)$; $\underline v\,(jm\omega + h +
k/j\omega) = B\ell\,\underline i$; substitua e resolva para $\underline v$.

\textbf{13.} $i \approx (u - B\ell v)/R$ dá $m\dot v + (h + B^2\ell^2/R)v + kx =
B\ell u/R$: amortecimento extra $25/8 = \qty{3.1}{N.s/m}$. $f_0 = \tfrac1{2\pi}\sqrt{k/m} =
\qty{50}{Hz}$; $Q = \sqrt{km}/h_{\text{tot}}$: $3.2$ em aberto, $0.77$ com o
amplificador.

\textbf{14.} $F = \qty{5.0}{N}$, $a = F/m = \qty{500}{m/s^2}$, $x = a/\omega^2 =
\qty{13}{\micro m}$, $v = a/\omega = \qty{0.080}{m/s}$; $B\ell v = \qty{0.40}{V}$ contra
$Ri = \qty{8}{V}$: $5\%$.

\textbf{15.} Com a corrente imposta, $v = F/h = \qty{5}{m/s}$ (uma
\omterm{def:b1:wave-propagation:signal}{amplitude} de \qty{16}{mm} --- uma excursão enorme para \qty{1}{A}); $B\ell v =
\qty{25}{V}$, três vezes $Ri$: o amplificador tem de fornecer \qty{33}{V} para empurrar
\qty{1}{A} na ressonância --- é o pico de \omterm{def:b1:sinusoidal-impedance:impedance}{impedância}. Um amplificador que fosse \omterm{def:b1:dc-circuits:sources}{fonte de tensão}
deixaria, em vez disso, o amortecimento elétrico segurar o cone.

\textbf{16.} $\tfrac12RI^2 = \qty{4.0}{W}$; $\tfrac12hv^2 = \tfrac12 \times 0.0064 =
\qty{3.2}{mW}$: \omterm{def:b1:heat-engines:efficiency}{rendimento} de $0.08\%$ --- os alto-falantes são aquecedores que
sussurram.

\textbf{17.} Acima da ressonância, $a = F/m = B\ell i/m$ não depende da
frequência para uma dada corrente, de modo que a pressão, $\propto a$, fica plana:
a banda controlada pela massa é a banda de trabalho.

\textbf{18.} Campo radial: a força é axial e a mesma em cada
posição da bobina no entreferro. Um tweeter tem de permanecer controlado pela massa
até \qty{20}{kHz} com uma excursão minúscula: bobina e cone leves; e um
cone pequeno irradia as frequências altas em todas as direções.

\textbf{19.} Baixa frequência: $v$ pequeno, $|Z| \approx R = \qty{8}{\ohm}$; a
\qty{50}{Hz}: $u = Ri + B\ell v$ com $v = B\ell i/h$: $Z = R + B^2\ell^2/h = 8 + 25 =
\qty{33}{\ohm}$; a \qty{10}{kHz}: $|R + jL\omega| = |8 + 31j| = \qty{32}{\ohm}$. Um
pico na ressonância, um piso plano de \qty{8}{\ohm} e uma subida com $L\omega$.

\textbf{20.} $e = B\ell v = 5 \times 10^{-3} = \qty{5}{mV}$.

\textbf{21.} $i = 5 \times 10^{-3}/608 = \qty{8.2}{\micro A}$; $F = B\ell i = \qty{41}{\micro N}$,
opondo-se a $v$: $F = [B^2\ell^2/(R + R_L)]\,v$, um amortecimento de \qty{0.041}{N.s/m};
a potência $Fv = \qty{4e-8}{W}$ é tirada do som e acaba nos
dois \omterm{def:b1:dc-circuits:resistor}{resistores}.

\textbf{22.} Carretel: $B^2\ell_f^2/R_f = 1 \times 0.01/10^{-3} = \qty{10}{N.s/m}$, dez
vezes $h$: ele amorteceria o cone e desperdiçaria potência como calor; uma fenda
rompe a espira fechada.

\textbf{23.} Em aberto: só $h$ amortece, $Q = 3.2$, e ele soa. Em curto: a
força contraeletromotriz faz passar corrente por $R$, acrescenta-se o amortecimento $B^2\ell^2/R$,
$Q = 0.77$: ele para na hora. (Experimente.)

\textbf{24.} Ao longo de um período, as energias armazenadas voltam aos seus valores:
$\langle ui\rangle = \langle Ri^2\rangle + \langle hv^2\rangle$ --- entrada elétrica $=$
calor Joule na bobina $+$ perdas mecânicas (som irradiado e
atrito da suspensão); o termo de acoplamento $B\ell iv$ aparece nas duas
equações com sinais opostos e se cancela.

\textbf{25.} Dínamo: Faraday numa bobina girante --- \qty{13.5}{V} a
\qty{18}{km/h}, com a corrente limitada a \qty{1.35}{A} pela sua própria indutância.
Alto-falante: \omterm{thm:b1:magnetostatics:laplace}{força de Laplace} e força contraeletromotriz --- \qty{3.1}{N.s/m} de amortecimento
elétrico, \omterm{def:b1:heat-engines:efficiency}{rendimento} de $0.08\%$. Microfone: a mesma bobina ao contrário ---
\qty{5}{mV} por milímetro por segundo.
\end{solution}
